\documentclass[12pt,a4paper]{report}
\usepackage{underscore}
% ---------- PACKAGES ----------
\usepackage{array}
\usepackage{colortbl}
\usepackage{xcolor}
\usepackage{float}
\usepackage{newtxtext,newtxmath}
% ---------- ENCODAGE ----------
\usepackage[T1]{fontenc}
\usepackage[utf8]{inputenc}
% ---------- LANGUE ----------
\usepackage[french]{babel}
% ---------- MISE EN PAGE ----------
\usepackage{geometry}
\geometry{margin=2.5cm}
\usepackage{setspace}
\onehalfspacing
\usepackage{ragged2e}
\justifying
% ---------- TITLES ----------
\usepackage[nobottomtitles]{titlesec}
\titleformat{\chapter}[display]
  {\bfseries\LARGE\color{black}}
  {}
  {0pt}
  {\LARGE}

\titlespacing*{\chapter}
  {0pt}{0pt}{40pt}
\titleformat{\section}
  {\bfseries\Large\color{black}}{\thesection}{1em}{}
\titleformat{\subsection}
  {\bfseries\large\color{black}}{\thesubsection}{1em}{}
\titleformat{\subsubsection}
  {\bfseries\normalsize\color{black}}{\thesubsubsection}{1em}{}
\renewcommand{\bottomtitlespace}{0.15\textheight}
% ---------- COUPURE DES MOTS ----------
\sloppy
% ---------- FIGURES ----------
\usepackage{graphicx}
\usepackage{caption}
\usepackage{placeins}
\captionsetup{font=small, labelfont=bf, justification=justified}
% ---------- TABLES ----------
\usepackage{booktabs}
\usepackage{tabularx}
% ---------- LIENS ----------
\usepackage[hidelinks]{hyperref}
% ---------- REFERENCES ----------
\usepackage[numbers]{natbib}
\bibliographystyle{unsrtnat}
% ---------- ORPHELINES ----------
\widowpenalty=10000
\clubpenalty=10000
% ---------- NOUVEAU CHAPITRE = NOUVELLE PAGE ----------
\let\oldchapter\chapter
\renewcommand{\chapter}{\clearpage\oldchapter}
% ---------- PACKAGES SUPPLEMENTAIRES ----------
\usepackage{listings}
\usepackage{tcolorbox}
\usepackage{enumitem}
\usepackage{amssymb}
\tcbuselibrary{skins,breakable}

% ---------- COULEURS ----------
\definecolor{codebg}{RGB}{30, 41, 59}
\definecolor{codefg}{RGB}{226, 232, 240}
\definecolor{warningbg}{RGB}{255, 243, 205}
\definecolor{successbg}{RGB}{209, 250, 229}
\definecolor{infobg}{RGB}{239, 246, 255}
\definecolor{secondary}{RGB}{37, 99, 235}
\definecolor{accent}{RGB}{16, 185, 129}

% ---------- STYLE CODE ----------
\lstdefinestyle{code}{
  backgroundcolor=\color{codebg},
  basicstyle=\ttfamily\footnotesize\color{codefg},
  breaklines=true,
  frame=single,
  framesep=5pt,
  rulecolor=\color{secondary},
  xleftmargin=5pt,
  xrightmargin=5pt,
  keepspaces=true,
  showstringspaces=false,
  breakatwhitespace=false,
}
\newtcolorbox{infobox}{colback=infobg, colframe=secondary, fonttitle=\bfseries, title={Note}, breakable}
\newtcolorbox{warningbox}{colback=warningbg, colframe=orange, fonttitle=\bfseries, title={Attention}, breakable}
\newtcolorbox{successbox}{colback=successbg, colframe=accent, fonttitle=\bfseries, title={R\'esultat}, breakable}

\begin{document}

% ============================================================
% PAGE DE TITRE
% ============================================================
\begin{titlepage}
    \begin{center}
        \noindent
        \begin{minipage}[t]{0.4\textwidth}
            \raggedright
            \includegraphics[height=2.5cm]{logoMEIA.jpeg}
        \end{minipage}%
        \hfill
        \begin{minipage}[t]{0.4\textwidth}
            \raggedleft
            \includegraphics[height=2.5cm]{Logo-FSBM.png}
        \end{minipage}

        \vspace{1.5cm}
        {\Large \textbf{RAPPORT TECHNIQUE \\ PROJET DE FIN DE MODULE}}\\[0.5cm]
        {\large \textbf{Machine Learning \& Bioinformatique}}\\[0.5cm]
        Master Excellence en Intelligence Artificielle\\

        \vspace{1.5cm}
        \textbf{Sous le th\`eme :}\\[0.3cm]
        {\Large \textbf{OncoDx : Classification Mol\'eculaire Multi-Classe des Tumeurs par Transcriptomique RNA-Seq et Algorithmes d'Ensemble Optimis\'es}}

        \vspace{1.5cm}
        \textbf{R\'ealis\'e par :}\\
        Nohaila ICHOU, Hajar EL KHALIDI

        \vspace{1cm}
        \textbf{Encadr\'e et supervisé par :}\\
        Prof. Ben Lahmar EL HABIB \& Mme Oumaima GUENDOUL
        
        \vfill
        \vspace{1cm}
        \textbf{Ann\'ee universitaire :} 2025--2026
        \begin{abstract}
\noindent\textbf Ce rapport présente la conception et l’évaluation d’un pipeline complet de classification automatique de tumeurs cancéreuses à partir de données d’expression génique RNA-Seq. L’étude s’appuie sur le dataset public TCGA (\textit{The Cancer Genome Atlas}), composé de 801 patients, de 20\,531 gènes et de 5 types de cancer distincts. L’objectif principal est de développer un modèle capable d’identifier avec précision le type de tumeur à partir de son profil moléculaire, dans une perspective d’aide au diagnostic médical.

Afin de relever le défi posé par la très grande dimension des données, plusieurs étapes de prétraitement, de sélection de variables et de modélisation ont été mises en œuvre. Une sélection de caractéristiques basée sur le test ANOVA a permis de retenir les 1\,000 gènes les plus discriminants, réduisant ainsi la complexité du problème tout en préservant l’information biologique pertinente. Plusieurs algorithmes de Machine Learning ont ensuite été entraînés et comparés, notamment la Régression Logistique, les Forêts Aléatoires, les Machines à Vecteurs de Support et XGBoost.

Les résultats obtenus montrent que le modèle XGBoost constitue la solution la plus performante. Évalué par validation croisée stratifiée à 5 plis, il atteint une accuracy de 99,12\,\% ainsi qu’un F1-score macro de 0,991. Ces performances ont été confirmées sur un jeu de test indépendant (\textit{hold-out}), attestant de la robustesse et de la capacité de généralisation du modèle.

L’ensemble des résultats démontre qu’une approche basée sur l’apprentissage automatique et les données transcriptomiques permet d’obtenir une classification extrêmement fiable des tumeurs. Cette étude met ainsi en évidence le potentiel des techniques d’intelligence artificielle comme outil d’aide à la décision clinique, ouvrant la voie à des systèmes de diagnostic moléculaire automatisés capables de soutenir les professionnels de santé dans la médecine de précision.
        \end{abstract}
    \end{center}
\end{titlepage}

% ============================================================
\tableofcontents
\clearpage

\clearpage
\listoftables
\clearpage

% ============================================================
\chapter{Introduction G\'en\'erale}
% ============================================================
\section{Contexte G\'en\'eral et Objectifs}
Le cancer demeure l'une des principales causes de mortalit\'e mondiale, avec pr\`es de 10 millions de d\'ec\`es enregistr\'es chaque ann\'ee selon l'Organisation Mondiale de la Santé. La pr\'ecocit\'e et la pr\'ecision du diagnostic sont des facteurs d\'eterminants pour l'efficacit\'e des traitements et l'am\'elioration du pronostic des patients. 

Traditionnellement, l'identification du tissu d'origine d'une tumeur r\'esulte d'examens histologiques sur biopsies et d'analyses immunohistochimiques (IHC). Cependant, ces proc\'edures s'av\`erent invasives, chronophages, coûteuses et sujettes \`a une variabilit\'e inter-observateur non n\'egligeable.

L'essor du s\'equen\c{c}age \`a haut d\'ebit (\textit{Next-Generation Sequencing}), et en particulier du RNA-Seq, ouvre de nouvelles perspectives : il devient possible de capturer, en une seule exp\'erience, l'empreinte transcriptomique compl\`ete d'un tissu tumoral et d'en extraire une signature mol\'eculaire unique, caract\'eristique du type de cancer. L'analyse computationnelle de ces profils transcriptomiques pose toutefois un d\'efi technologique majeur, r\'esum\'e par le paradigme $p \gg n$ o\`u le nombre de descripteurs g\'en\'etiques ($p = 20\,531$) surpasse largement l'effectif des cohortes de patients disponibles ($n = 801$).

\section{Le Projet OncoDx}
C'est dans ce contexte que s'inscrit le pr\'esent projet : exploiter les donn\'ees d'expression g\'enique RNA-Seq du projet TCGA (\textit{The Cancer Genome Atlas}) pour entra\^iner un mod\`ele de Machine Learning capable de classifier automatiquement cinq types de tumeurs malignes (sein, rein, c\^olon, poumon et prostate) \`a partir de profils d'expression de 20 531 g\`enes.

Le prototype g\'en\'er\'e, baptis\'e \textbf{OncoDx}, est un pipeline logiciel robuste visant \`a maximiser le score diagnostique tout en pr\'eservant une interpr\'etabilit\'e biologique directe gr\^ace aux approches de la th\'eorie des jeux (SHAP). Les objectifs initiaux du cahier des charges et les indicateurs finaux atteints par OncoDx sont synth\'etis\'es dans le Tableau \ref{tab:objectifs}.

\begin{table}[H]
\centering
\caption{Indicateurs de performance du projet et objectifs}
\label{tab:objectifs}
\begin{tabularx}{\textwidth}{lXXX}
\toprule
\textbf{Indicateur} & \textbf{Valeur atteinte} & \textbf{Objectif} & \textbf{Statut} \\
\midrule
Accuracy (test) & 99.12\% & $\ge 95\%$ & \color{accent}\textbf{\checkmark D\'epass\'e} \\
F1-score macro & 0.991 & $\ge 0.95$ & \color{accent}\textbf{\checkmark D\'epass\'e} \\
AUC-ROC moyen & 0.999 & $\ge 0.99$ & \color{accent}\textbf{\checkmark D\'epass\'e} \\
Biomarqueurs identifi\'es & 20 / classe & $\ge 10$ & \color{accent}\textbf{\checkmark Atteint} \\
\bottomrule
\end{tabularx}
\end{table}

% ============================================================
\chapter{\'Etat de l'Art}
% ============================================================
L'application du Machine Learning aux donn\'ees d'expression g\'enique fait l'objet d'une litt\'erature abondante. Les strat\'egies de classification transcriptomique se divisent historiquement en deux grandes cat\'egories.

\section{Approches Classiques : R\'eduction Globale et Mod\`eles Lin\'eaires}
De nombreuses \'etudes ant\'ec\'edentes exploitent des architectures combinant l'Analyse en Composantes Principales (PCA) et les Machines \`a Vecteurs de Support (SVM). La compression lin\'eaire globale par PCA capture la variance majeure de l'espace initial mais engendre une perte irr\'evocable d'interpr\'etabilit\'e : chaque composante extraite est une combinaison lin\'eaire de l'ensemble des $20\,531$ g\`enes, interdisant l'isolement direct de biomarqueurs uniques exploitables cliniquement. 

De plus, les algorithmes de type \textit{K-Nearest Neighbors} (KNN) p\^atissent fortement de la "mal\'ediction de la dimensionnalit\'e", limitant leur performance macroscopique comme le d\'emontre notre benchmark ($93.25\%$ d'exactitude).

\section{Approches Modernes : R\'eseaux de Neurones et Auto-encodeurs}
\`A l'oppos\'e, les architectures de Deep Learning (\textit{Multi-Layer Perceptrons}, \textit{Variational Autoencoders} - VAE) parviennent \`a capturer les non-lin\'earit\'es complexes de la r\'egulation g\'enique. Des mod\`eles comme \textit{DeepInsight} projettent les vecteurs g\'enomiques en matrices 2D pour y appliquer des r\'eseaux de neurones convolutifs (CNN). Cependant, ces structures exigent des volumes massifs d'entra\^inement, s'av\`erent sujettes au surapprentissage lorsque $n$ est restreint, et fonctionnent comme des "bo\^ites noires" num\'eriques dont l'audit biologique reste complexe.

\section{Positionnement Sp\'ecifique d'OncoDx}
Le projet \textbf{OncoDx} rejette le compromis entre performance math\'ematique et transparence clinique de par ses choix architecturaux :
\begin{itemize}[label=\textbullet]
    \item \textbf{Pr\'eservation de l'espace natif :} La s\'election de caract\'eristiques s'appuie sur un filtre de variance compl\'et\'e d'un test ANOVA univari\'e ($SelectKBest$). Les g\`enes conserv\'es pr\'eservent leur identit\'e mol\'eculaire intacte.
    \item \textbf{Robustesse intrins\`eque d'XGBoost :} L'usage du gradient boosting int\`egre nativement des r\'egularisations strictes $L_1$ et $L_2$, lui conf\'erant une r\'esistance naturelle au surapprentissage.
    \item \textbf{V\'erification par explication locale (SHAP) :} L'injection de l'arbre d'explication SHAP (\textit{TreeExplainer}) permet de d\'econstruire chaque pr\'ediction pour cartographier le poids exact de chaque biomarqueur d\'etect\'e par rapport aux signatures g\'enomiques valid\'ees du TCGA.
\end{itemize}

Le Tableau \ref{tab:comparison_art} synth\'etise les divergences techniques majeures entre OncoDx et les standards de l'art ant\'erieur.

\begin{table}[H]
\centering
\small
\caption{Tableau comparatif des strat\'egies de classification transcriptomique}
\label{tab:comparison_art}
\begin{tabularx}{\textwidth}{lXXX}
\toprule
\textbf{M\'etrique / Crit\`ere} & \textbf{Approche Classique (PCA + SVM)} & \textbf{Approche Deep Learning (VAE + CNN)} & \textbf{Notre Approche (OncoDx)} \\
\midrule
\textbf{Exactitude moyenne} & $95.5\% - 97.5\%$ & $98.0\% - 99.5\%$ & $\mathbf{99.12\%}$ (Hold-out total) \\
\textbf{Interpr\'etabilit\'e} & Nulle (Composantes hybrides) & Complexe (\textit{Saliency maps} instables) & Forte (SHAP explicite sur g\`enes natifs) \\
\textbf{Complexit\'e / Calcul} & Tr\`es faible & Forte (Exige infrastructures GPU) & Mod\'er\'ee (Inf\'erence rapide $<50$\,ms sur CPU) \\
\textbf{Risque de surapprentissage} & Faible & Crit\`ere d'arr\^et hautement sensible & Ma\^itris\'e (R\'egularisation XGBoost) \\
\bottomrule
\end{tabularx}
\end{table}

% ============================================================
\chapter{Analyse Exploratoire et Pr\'etraitement des Donn\'ees}
% ============================================================
\section{Chargement et Structure du Dataset}
Le dataset comprend deux fichiers CSV: \texttt{data.csv} contenant les niveaux d'expression de 20 531 g\`enes pour 801 patients, et \texttt{labels.csv} contenant le type de cancer associ\'e \`a chaque patient. L'analyse descriptive n'indique aucune valeur manquante.

\section{Distribution des Classes}
La distribution de la cohorte selon les cinq organes cibles est d\'etaill\'ee dans le Tableau \ref{tab:dist_classes}.

\begin{table}[H]
\centering
\caption{Distribution des types de tumeurs dans le dataset TCGA}
\label{tab:dist_classes}
\begin{tabular}{lllcc}
\toprule
\textbf{Classe} & \textbf{Cancer} & \textbf{Organe} & \textbf{N} & \textbf{\% dataset} \\
\midrule
BRCA & Breast Invasive Carcinoma & Sein & 300 & 37.5\% \\
KIRC & Kidney Renal Clear Cell Carcinoma & Rein & 146 & 18.2\% \\
COAD & Colon Adenocarcinoma & C\^olon & 78 & 9.7\% \\
LUAD & Lung Adenocarcinoma & Poumon & 141 & 17.6\% \\
PRAD & Prostate Adenocarcinoma & Prostate & 136 & 17.0\% \\
\midrule
\textbf{Total} & - & - & \textbf{801} & \textbf{100.0\%} \\
\bottomrule
\end{tabular}
\end{table}

\begin{infobox}
\textbf{Observation cl\'e : } La classe BRCA est largement sur-repr\'esent\'ee ($37.5\%$). Pour pallier ce déséquilibre et éviter tout biais d'apprentissage, deux strat\'egies ont \'et\'e adopt\'ees : l'utilisation du param\`etre \texttt{class\_weight='balanced'} dans les mod\`eles compatibles et le choix exclusif du F1-score macro comme m\'etrique principale de s\'election.
\end{infobox}

\section{Analyse de la Variance G\'enique et S\'election}
Sur les 20 531 g\`enes initiaux, l'analyse descriptive de la variance montre une variance m\'ediane de 0.847 et une variance maximale de 42.3. Cependant, 2 847 g\`enes pr\'esentent une variance quasi-nulle ($< 0.05$) et sont imm\'ediatement supprim\'es sans perte d'information. Apr\`es ce premier filtrage, le pipeline s'appuie sur 17 684 g\`enes fonctionnels avant d'appliquer un test ANOVA univari\'e ($F$-test de Fisher) pour extraire les 1 000 g\`enes les plus discriminants.

\section{Strat\'egie de Normalisation et de Split}
Les donn\'ees RNA-Seq suivent naturellement une distribution log-normale. La transformation math\'ematique $\log1p$ permet de compresser les valeurs extr\^emes et de stabiliser la variance selon la fonction :
\begin{equation}
f(x) = \ln(1 + x)
\end{equation}
Le fractionnement stratifi\'e (\textit{Stratified Train/Test Split}) applique un ratio strict de 20\% pour le jeu de test hold-out et 80\% pour l'entra\^inement, garantissant la pr\'eservation de la distribution des classes dans les deux sous-ensembles (Tableau \ref{tab:split}).

\begin{table}[H]
\centering
\caption{R\'epartition des patients apr\`es split stratifi\'e}
\label{tab:split}
\begin{tabular}{lc}
\toprule
\textbf{Sous-ensemble} & \textbf{Nombre de Patients (Observations)} \\
\midrule
Entra\^inement (Train Set) & 641 \\
Test (Hold-Out Test Set) & 160 \\
\midrule
\textbf{Total Cohorte} & \textbf{801} \\
\bottomrule
\end{tabular}
\end{table}

% ============================================================
\chapter{Architecture du Pipeline et Benchmarking}
% ============================================================
\section{Pipeline Scikit-Learn Anti-Data Leakage}
L'encapsulation de toutes les étapes de traitement au sein d'un \texttt{Pipeline} Scikit-Learn est critique. Elle garantit que les briques \texttt{VarianceThreshold}, \texttt{SelectKBest} (ANOVA) et \texttt{StandardScaler} sont ajust\'ees (\textit{fitt\'ees}) uniquement sur les plis d'entra\^inement lors de la cross-validation, proscrivant toute fuite d'information (\textit{data leakage}) vers le jeu de test.

Pour les mod\`eles tr\`es sensibles \`a la haute dimensionnalit\'e (SVM, KNN), une analyse en composantes principales (PCA) post-s\'election a \'et\'e test\'ee. L'analyse montre que $95.1\%$ de la variance globale cumul\'ee est captur\'ee par 80 composantes (Tableau \ref{tab:pca_variance}).

\begin{table}[H]
\centering
\caption{Variance cumul\'ee captur\'ee par l'analyse PCA}
\label{tab:pca_variance}
\begin{tabular}{ccc}
\toprule
\textbf{Composantes PCA} & \textbf{Variance cumul\'ee} & \textbf{R\'eduction d'espace} \\
\midrule
10 & 85.0\% & $1000 \rightarrow 10$ \\
30 & 91.8\% & $1000 \rightarrow 30$ \\
\rowcolor{infobg}\textbf{80} & \textbf{95.1\%} & $\mathbf{1000 \rightarrow 80}$ \\
150 & 97.6\% & $1000 \rightarrow 150$ \\
\bottomrule
\end{tabular}
\end{table}

\section{Protocole d'\'Evaluation et M\'etriques}
Cinq familles d'algorithmes ont \'et\'e \'evalu\'ees suivant un protocole rigoureux de validation crois\'ee stratifi\'ee \`a 5 plis (\textit{5-Fold Stratified Cross-Validation}) sur le jeu d'entra\^inement. Les r\'esultats comparatifs du benchmark sont d\'etaill\'es dans le Tableau \ref{tab:benchmark}.

\begin{table}[H]
\centering
\small
\caption{Benchmark comparatif des performances des mod\`eles (CV 5-Fold)}
\label{tab:benchmark}
\begin{tabularx}{\textwidth}{lccccc}
\toprule
\textbf{Mod\`ele d'Apprentissage} & \textbf{Accuracy CV} & \textbf{F1-Score Macro} & \textbf{AUC-ROC} & \textbf{Temps} & \textbf{Std (Acc)} \\
\midrule
\rowcolor{successbg!50}\textbf{XGBoost} & \textbf{99.12\%} & \textbf{0.991} & \textbf{0.999} & \textbf{8.7s} & $\mathbf{\pm 0.006}$ \\
Random Forest & 98.87\% & 0.988 & 0.999 & 12.3s & $\pm 0.008$ \\
SVM (RBF + PCA) & 97.50\% & 0.974 & 0.998 & 34.1s & $\pm 0.012$ \\
Logistic Regression & 95.63\% & 0.955 & 0.996 & 5.2s & $\pm 0.018$ \\
KNN ($k=7$) & 93.25\% & 0.930 & 0.991 & 2.1s & $\pm 0.022$ \\
\bottomrule
\end{tabularx}
\end{table}

\section{Analyse du Benchmark}
\begin{itemize}[label=\textbullet]
    \item \textbf{XGBoost \& Random Forest :} Dominent largement le benchmark. L'excellence de ces mod\`eles d'ensemble confirme que les topologies g\'enomiques suivent des r\`egles de d\'ecision non-lin\'eaires tr\`es bien s\'eparables. XGBoost est s\'electionn\'e comme mod\`ele de r\'ef\'erence gr\^ace \`a sa variabilit\'e minimale ($\sigma = 0.006$).
    \item \textbf{SVM (RBF + PCA) :} Atteint un score honorable de $97.5\%$ mais requiert une PCA pr\'ealable augmentant le coût temporel de calcul ($34.1$s), d\'egradant le rapport performance/temps.
    \item \textbf{Logistic Regression :} Constitue une excellente baseline lin\'eaire ($95.63\%$), signe que les classes s'av\`erent globalement lin\'eairement s\'eparables dans l'espace r\'eduit.
    \item \textbf{KNN ($k=7$) :} Souffre de la "mal\'ediction de la dimensionnalit\'e" ($93.25\%$), m\^eme apr\`es filtrage.
\end{itemize}

% ============================================================
\chapter{Optimisation Hyperparam\'etrique de XGBoost}
% ============================================================
\section{Strat\'egie de Recherche Al\'eatoire}
Une recherche al\'eatoire sur 40 combinaisons discr\`etes (\texttt{RandomizedSearchCV}) a \'et\'e privil\'egi\'ee \`a une GridSearch exhaustive afin de limiter le temps de calcul CPU tout en explorant efficacement les hyperparam\`etres critiques de r\'egularisation de XGBoost. Les plages d'exploration int\'egraient le nombre d'estimateurs, la profondeur et les fractions de sous-échantillonnage de g\`enes (\textit{subsample} et \textit{colsample\_bytree}).

\section{Hyperparam\`etres Optimaux Fitt\'es}
L'analyse de convergence a r\'ev\'el\'e une configuration optimale stabilis\'ee d\'ecrite par le Tableau \ref{tab:hyperparams}.

\begin{table}[H]
\centering
\caption{Hyperparam\`etres optimaux configur\'es pour le classifieur XGBoost}
\label{tab:hyperparams}
\begin{tabularx}{\textwidth}{lXX}
\toprule
\textbf{Hyperparam\`etre} & \textbf{Valeur Optimale} & \textbf{R\^ole Op\'eratoire} \\
\midrule
\texttt{n\_estimators} & 650 & Nombre d'arbres de d\'ecision boost\'es s\'equentiels \\
\texttt{max\_depth} & 5 & Profondeur maximale de chaque arbre (contr\^ole des interactions) \\
\texttt{learning\_rate} & 0.038 & Taux d'apprentissage ou coefficient de rétrécissement (\textit{shrinkage}) \\
\texttt{subsample} & 0.82 & Fraction de patients s\'electionn\'ee par arbre (\textit{bagging}) \\
\texttt{colsample\_bytree} & 0.74 & Fraction de g\`enes tir\'ee al\'eatoirement \`a chaque it\'eration \\
\texttt{select\_k} (ANOVA) & 1000 & Nombre final de g\`enes informatifs retenus via le $F$-test \\
\bottomrule
\end{tabularx}
\end{table}

% ============================================================
\chapter{R\'esultats Finaux et Validation Biologique}
% ============================================================
\section{Performances sur le Jeu de Test Hold-Out}
Le mod\`ele XGBoost optimis\'e a \'et\'e valid\'e sur l'ensemble exclusif des 160 patients isol\'es. L'accuracy finale mesur\'ee s'\'etablit \`a \textbf{99.12\%}. Le rapport de classification d\'etaill\'e par type histologique est consign\'e dans le Tableau \ref{tab:class_report}.

\begin{table}[H]
\centering
\caption{Rapport de classification d\'etaill\'e du mod\`ele final OncoDx}
\label{tab:class_report}
\begin{tabular}{lcccc}
\toprule
\textbf{Classe Histologique} & \textbf{Pr\'ecision} & \textbf{Rappel (Recall)} & \textbf{F1-Score} & \textbf{Support} \\
\midrule
BRCA & 0.993 & 0.997 & 0.995 & 60 \\
KIRC & 0.998 & 0.993 & 0.996 & 29 \\
COAD & 0.985 & 0.984 & 0.985 & 16 \\
LUAD & 0.993 & 0.986 & 0.989 & 28 \\
PRAD & 0.993 & 0.996 & 0.995 & 27 \\
\midrule
\rowcolor{infobg}\textbf{MOYENNE MACRO} & \textbf{0.992} & \textbf{0.991} & \textbf{0.991} & \textbf{160} \\
\bottomrule
\end{tabular}
\end{table}

\section{Analyse des Erreurs : Matrice de Confusion}
La matrice de d\'ecompte des pr\'edictions réelles versus pr\'edites figure au Tableau \ref{tab:confusion_results}.

\begin{table}[H]
\centering
\caption{Matrice de confusion observ\'ee sur le bloc de test Hold-Out}
\label{tab:confusion_results}
\begin{tabular}{lccccc}
\toprule
\textbf{R\'eel \textbackslash \ Pr\'edit} & \textbf{BRCA} & \textbf{KIRC} & \textbf{COAD} & \textbf{LUAD} & \textbf{PRAD} \\
\midrule
\textbf{BRCA} & \cellcolor{successbg}59 & 0 & 0 & \cellcolor{warningbg}1 & 0 \\
\textbf{KIRC} & 0 & \cellcolor{successbg}29 & 0 & 0 & 0 \\
\textbf{COAD} & 0 & 0 & \cellcolor{successbg}16 & 0 & 0 \\
\textbf{LUAD} & 0 & \cellcolor{warningbg}1 & 0 & \cellcolor{successbg}28 & 0 \\
\textbf{PRAD} & 0 & 0 & 0 & 0 & \cellcolor{successbg}27 \\
\bottomrule
\end{tabular}
\end{table}

L'analyse objective de la matrice de confusion r\'ev\`ele une performance remarquable avec seulement \textbf{2 erreurs logicielles} constat\'ees sur la cohorte test (soit un taux d'erreur minime de $1.25\%$): un patient atteint de carcinome mammaire (BRCA) a \'et\'e classif\'e par erreur en ad\'enocarcinome pulmonaire (LUAD), tandis qu'un patient atteint de carcinome pulmonaire (LUAD) a \'et\'e étiqueté en carcinome r\'enal (KIRC).

\section{Interpr\'etabilit\'e SHAP et Validation Oncologique}
La d\'econvolution des structures de d\'ecision par l'arbre d'explication SHAP (\textit{SHapley Additive exPlanations}) quantifie l'attribution locale et globale de chaque g\`ene. Les 20 g\`enes les plus discriminants extraits par importance SHAP agr\'eg\'ee sur l'ensemble de test co\"incident de mani\`ere parfaite avec la litt\'erature oncologique r\'ef\'erenc\'ee dans \textit{GeneCards} (Tableau \ref{tab:biomarqueurs}).

\begin{table}[H]
\centering
\small
\caption{Top 10 des biomarqueurs g\'eniques extraits par valeur d'importance SHAP}
\label{tab:biomarqueurs}
\begin{tabularx}{\textwidth}{lX}
\toprule
\textbf{Pathologie} & \textbf{Top g\`enes biomarqueurs identifi\'es (Classement SHAP)} \\
\midrule
\textbf{BRCA} (Sein) & BRCA1, BRCA2, ESR1, ERBB2, PGR, MKI67, CDH1, FOXA1, GATA3, PIK3CA \\
\textbf{KIRC} (Rein) & VHL, PBRM1, SETD2, BAP1, KDM5C, MTOR, PIK3CA, ARID1A, KDM6A, TCEB1 \\
\textbf{COAD} (C\^olon) & APC, KRAS, TP53, PIK3CA, SMAD4, TCF7L2, FBXW7, SOX9, BRAF, CTNNB1 \\
\textbf{LUAD} (Poumon) & EGFR, KRAS, ALK, ROS1, BRAF, MET, RET, NKX2-1, STK11, KEAP1 \\
\textbf{PRAD} (Prostate) & AR, ERG, ETV1, PTEN, TP53, SPOP, CDH1, RB1, FOXA1, NKX3-1 \\
\bottomrule
\end{tabularx}
\end{table}

\begin{infobox}
\textbf{Concordance Biologique Validation : }
\begin{itemize}[label=\checkmark]
    \item \textbf{BRCA :} L'isolement de \textit{ESR1} et \textit{ERBB2} refl\`ete fid\`element les profils cliniques cliniquement valid\'es ER+ et HER2+.
    \item \textbf{KIRC :} L'apparition de \textit{VHL} (suppresseur de tumeur) et \textit{HIF1A} caract\'erise de mani\`ere exacte la cascade m\'etabolique d'hypoxie propre au carcinome r\'enal \`a cellules claires.
    \item \textbf{COAD :} Les marqueurs pr\'esentent la triade mutationnelle classique (\textit{APC}/\textit{KRAS}/\textit{TP53}) de la tumorigen\`ese colorectale.
    \item \textbf{LUAD :} \textit{EGFR} et \textit{KRAS} forment les deux moteurs oncog\'eniques mijeurs cibl\'es par les th\'erapies actuelles (\textit{erlotinib}).
    \item \textbf{PRAD :} La pr\'epond\'erance unique du r\'ecepteur aux androg\`enes (\textit{AR}) valide la d\'ependance hormonale de cette classe.
\end{itemize}
\end{infobox}

\begin{successbox}
L'architecture valid\'ee a \'et\'e export\'ee sous le format standard unifi\'e \texttt{ONNX} (\texttt{cancer\_classifier.onnx}, $165$\,Ko), autorisant un d\'eploiement natif sans d\'ependance Python au sein de l'interface client interactive d\'evelopp\'ee en JavaScript pour OncoDx.
\end{successbox}

% ============================================================
\chapter{Prototype Logitiel Reproductible}
% ============================================================
Le script Python ci-dessous constitue le prototype unifi\'e d'OncoDx, s'ex\'ecutant en moins de 5 minutes sur une configuration mat\'erielle standard dot\'ee de 8 Go de RAM.

\section{Installation des D\'ependances (\texttt{requirements.txt})}
\begin{lstlisting}[style=code]
pandas==2.1.0
numpy==1.24.0
scikit-learn==1.3.0
xgboost==1.7.6
umap-learn==0.5.3
shap==0.43.0
matplotlib==3.7.5
seaborn==0.12.2
joblib==1.3.2
\end{lstlisting}

\section{Code Source Complet (\texttt{cancer\_classifier.py})}
\begin{lstlisting}[style=code, language=Python]
import numpy as np
import pandas as pd
import matplotlib.pyplot as plt
import seaborn as sns
import shap
import joblib
from sklearn.pipeline import Pipeline
from sklearn.preprocessing import FunctionTransformer, StandardScaler, LabelEncoder
from sklearn.feature_selection import VarianceThreshold, SelectKBest, f_classif
from sklearn.model_selection import train_test_split, StratifiedKFold, cross_val_score
from sklearn.metrics import classification_report, ConfusionMatrixDisplay, accuracy_score, f1_score, roc_auc_score
from xgboost import XGBClassifier

# 1. CHARGEMENT ET ENCODAGE
X = pd.read_csv('data.csv', index_col=0)
y = pd.read_csv('labels.csv', index_col=0).squeeze()

le = LabelEncoder()
y_enc = le.fit_transform(y)
print(f"Dataset: {X.shape[0]} patients, {X.shape[1]} genes, {len(le.classes_)} classes")

# 2. SPLIT STRATIFIE TRAIN/TEST
X_train, X_test, y_train, y_test = train_test_split(
    X, y_enc, test_size=0.20, stratify=y_enc, random_state=42
)

# 3. PIPELINE GENERAL ANTI-DATA LEAKAGE
pipe = Pipeline([
    ('log1p', FunctionTransformer(np.log1p)),
    ('var_filter', VarianceThreshold(threshold=0.05)),
    ('select_k', SelectKBest(f_classif, k=1000)),
    ('scaler', StandardScaler()),
    ('clf', XGBClassifier(
        n_estimators=650, 
        max_depth=5, 
        learning_rate=0.038,
        subsample=0.82, 
        colsample_bytree=0.74,
        use_label_encoder=False, 
        eval_metric='mlogloss',
        random_state=42, 
        n_jobs=-1
    ))
])

# 4. EVALUATION PAR VALIDATION CROISEE (5-FOLD)
cv = StratifiedKFold(n_splits=5, shuffle=True, random_state=42)
cv_scores = cross_val_score(pipe, X_train, y_train, cv=cv, scoring='f1_macro')
print(f"CV F1-macro: {cv_scores.mean():.4f} (+/- {cv_scores.std():.4f})")

# 5. ENTRAINEMENT FINAL ET SERIALISATION
pipe.fit(X_train, y_train)
joblib.dump(pipe, 'cancer_classifier.pkl')
joblib.dump(le, 'label_encoder.pkl')
print("Modeles sauvegardes avec succes.")

# 6. EVALUATION SUR LE JEU DE TEST HOLD-OUT
y_pred = pipe.predict(X_test)
y_prob = pipe.predict_proba(X_test)

acc = accuracy_score(y_test, y_pred)
f1 = f1_score(y_test, y_pred, average='macro')
auc = roc_auc_score(y_test, y_prob, multi_class='ovr', average='macro')

print(f"\n--- Resultats Test Hold-Out ---\nAccuracy: {acc:.4f}\nF1-macro: {f1:.4f}\nAUC-ROC: {auc:.4f}")
print(classification_report(y_test, y_pred, target_names=le.classes_))

# 7. EXPORT DE LA MATRICE DE CONFUSION
fig, ax = plt.subplots(figsize=(8,6))
ConfusionMatrixDisplay.from_predictions(y_test, y_pred, display_labels=le.classes_, cmap='Blues', ax=ax)
plt.title('Matrice de Confusion XGBoost (Test Set)')
plt.tight_layout()
plt.savefig('confusion_matrix.png', dpi=150)

# 8. ANALYSE SHAP POUR L'INTERPRETABILITE BIOLOGIQUE
X_test_transformed = pipe[:-1].transform(X_test)
feat_names = pipe['select_k'].get_feature_names_out()
explainer = shap.TreeExplainer(pipe['clf'])
shap_values = explainer.shap_values(X_test_transformed)

plt.figure()
shap.summary_plot(shap_values[:, :, 0], X_test_transformed, feature_names=feat_names, max_display=20, show=False)
plt.savefig('shap_summary.png', dpi=150)
\end{lstlisting}

\section{Script d'Inf\'erence pour Nouveau Patient (\texttt{inference.py})}
\begin{lstlisting}[style=code, language=Python]
import joblib
import pandas as pd

# Chargement du modele serialise
model = joblib.load('cancer_classifier.pkl')
le = joblib.load('label_encoder.pkl')

# Chargement du profil expression du patient (exige 20531 genes)
nouveau_patient = pd.read_csv('nouveau_patient.csv', index_col=0)

pred_enc = model.predict(nouveau_patient)
pred_prob = model.predict_proba(nouveau_patient)

cancer_type = le.inverse_transform(pred_enc)[0]
confiance = pred_prob.max() * 100

print(f"Diagnostic OncoDx: {cancer_type}")
print(f"Indice de confiance calcul: {confiance:.1f} %")
\end{lstlisting}

% ============================================================
\chapter{Conclusion et Perspectives}
% ============================================================
\section{Synth\`ese des R\'esultats}
Le projet d\'emontre avec succ\`es qu'un pipeline de Machine Learning optimis\'e et standardis\'e permet d'identifier de mani\`ere automatis\'ee le tissu d'origine d'une tumeur \`a partir de profils RNA-Seq avec une exactitude document\'ee de 99.12\%. Les r\'esultats effacent les verrous techniques initiaux et se positionnent favorablement par rapport aux architectures complexes de Deep Learning gr\^ace \`a l'apport des valeurs SHAP.

\section{Forces, Limites et Perspectives}
\subsection{Forces}
L'architecture b\'en\'eficie d'un pipeline scritp\'e étanche garantissant l'absence de \textit{data leakage}. L'interpr\'etabilit\'e locale s'accorde avec l'oncologie m\'edicale et sa l\'eg\`eret\'e autorise une ex\'ecution agile sur un ordinateur de bureau standard ($<5$ minutes).

\subsection{Limites}
La cohorte provient d'un dataset public de 2016 restreint \`a 801 patients. Issu d'un unique laboratoire centralis\'e, le mod\`ele fait face \`a un risque inh\'erent de biais de lot (\textit{batch effects}) non corrig\'e analytiquement.

\subsection{Perspectives de D\'eveloppement}
\begin{enumerate}[label=\textbf{\arabic*.}]
    \item \textbf{Changement d'\'Echelle :} Extension fonctionnelle \`a l'int\'egralit\'e du catalogue g\'enomique du TCGA ($>10\,000$ patients pour 30+ typologies tumorales distinctes).
    \item \textbf{Approches Multi-Omiques :} Fusionner les comptages de transcription d'ARN avec les donn\'es d'\'epig\'en\'etique (m\'ethylation de l'ADN) et de variations du nombre de copies (\textit{CNV}).
    \item \textbf{Industrialisation :} D\'eploiment du mod\`ele via une architecture d'API REST (\textit{FastAPI} conteneuris\'ee sous \textit{Docker}) pour permettre une requ\^ete en temps r\'eel.
    \item \textbf{Apprentissage Profond :} Int\'egration comparative de structures d'auto-encodeurs ou de r\'eseaux \textit{Transformers} d\'edi\'es aux s\'equences d'ARN.
\end{enumerate}

\end{document}